\begin{center}\begin{tikzpicture}[font=\small,>=stealth,line join=round]
\draw[gray!40,rounded corners=3pt] (0,0) rectangle (14.2,8.3);\node at (7.1,7.700000000000001) {例4：偶函数的两支，给出两个可能的n};
\begin{scope}[shift={(4.1,2.161111111111111)},x=0.8571428571428571cm,y=1.0222222222222221cm]\draw[->,gray] (-3.5,0)--(3.5,0) node[right]{$x$};\draw[->,gray] (0,-0.5)--(0,4) node[above]{$y$};\draw[blue,thick,solid,domain=-3:-0.22,samples=180] plot(\x,{abs(\x)^(-2/3)});\draw[blue,thick,solid,domain=0.22:3,samples=180] plot(\x,{abs(\x)^(-2/3)});\fill[orange!85!black] (-0.5443310539518174,1.5) circle(2.2pt);\fill[orange!85!black] (0.5443310539518174,1.5) circle(2.2pt);\draw[orange!85!black,dashed] (-3,1.5)--(3,1.5);\node[left] at (-3,1.5){$\frac32$};\end{scope}\node[anchor=west,align=left,text width=5.0cm] at (8.4,4.0) {$f(x)=1/|x|^{2/3},\ x\ne0$\\[9pt]关于y轴对称\\[9pt]$f(n)=\frac32$\\[9pt]$|n|=\frac{2\sqrt6}{9}$\\[9pt]$n=\pm\frac{2\sqrt6}{9}$};
\node[text width=12.6cm,align=center] at (7.1,.65) {负数底数可以取奇次根；两支分别单调，不能把整个定义域当作一个递减区间。};
\end{tikzpicture}\end{center}
